\section{Introduction}\label{intro}

Open Source Software (OSS) is a cornerstone of modern software development, providing a collaborative environment that fosters innovation, knowledge exchange, and the widespread use of software tools and libraries. 
The licensing of OSS projects is crucial as it dictates how software can be used, modified, and distributed.
As the number of OSS projects has surged, understanding the intricacies of license selection has become increasingly complex. 

Existing work either lacks the scale or comprehensive identification of licenses to fully reflect license usage in OSS.
More importantly, licence choice, much like the choice of technology~\cite{choice23,ma2020methodology}, may reflect strategic decisions by project maintainers to balance openness, protection, and control over their codebases.
It may also be influenced by matters of technical convenience (e.g., licenses used by the frameworks they depend on), social convenience (e.g., licenses favored by their
collaborators), or simply by awareness, such as the popularity of certain licenses.
% **work backwards from the regresssion predictors to 
% see how they fit within the choice literature and within prior
% license work literature. For example, language and calendar time may be a proxy
% for GPL popularity with older products and C more exposed to  gpl***.

Prior work only tabulated observed license use, but have not attempted to 
% *check* to observe evolution of license type use over time, nor to
model explicitly how the context, such as project size, programming language and other aspects of the project (or individual developer) influence license choice. 

To address these gaps, we aim to develop a model that a) comprehensively accounts for all license types, including minor and non-semantic variations, b) is applied to a sample that approximates the full scale of the OSS ecosystem, and c) incorporates key contextual predictors that may influence differing priorities in license selection among projects and individuals.

The core questions our research seeks to answer are:
\begin{itemize}
    \item \textbf{RQ1}: What are the types, prevalence, and evolution of licenses across the entire open source software landscape?
    \item \textbf{RQ2}: How the characteristics of a project influence the choice of its license?   
\end{itemize}


% While previous research has often emphasized the popularity of
% permissive licenses such as the MIT or Apache licenses, known for
% their flexibility and minimal restrictions, this perspective does
% not fully capture the nuanced landscape of OSS licensing. 
% In particular, the roles of copyleft and weak copyleft licenses,
% which impose stricter conditions to maintain the open nature of
% derivative works, deserve closer examination. 

% To address these complexities, we adopt a data-driven approach,
% analyzing a comprehensive dataset that covers a wide array of OSS
% projects. 
% Our analysis is enhanced by advanced statistical models that account
% for various project characteristics, including the number of
% commits, project maturity, team composition, and file counts. 
% This methodology allows us to identify underlying trends and
% provides a deeper understanding of the factors influencing license
% selection. 


% By addressing these questions, we aim to provide valuable insights
% for developers, project maintainers, and the broader OSS community,
% helping them make informed licensing decisions. 
% These decisions are critical not only for legal compliance and risk
% mitigation but also for sustaining and growing open source
% projects. 

% Our study challenges the prevailing view that permissive licenses are universally favored in OSS by presenting evidence of a more varied licensing landscape.
% We demonstrate that copyleft licenses, often considered restrictive, continue to play a significant role, especially in older and more complex projects.
% Additionally, we emphasize the considerable legal risks posed by the large number of unlicensed projects, highlighting the need for greater awareness and education within the OSS community.

% By illuminating these dynamics, our research contributes to a more profound understanding of the factors that influence license selection in OSS.
% We believe these insights will be pivotal in guiding future research and policy development, ultimately strengthening the legal and operational foundations of the OSS ecosystem.
